\documentclass[../../../include/open-logic-section]{subfiles}

\begin{document}

\olfileid{lam}{syn}{abb}
\olsection{缩写语法}

\olref[trm]{defn:term} 所定义的项有时写起来颇为繁琐，因此引入一种更简洁的语法很有用。
当然，我们必须谨慎地确保，用这种简洁语法写出的项也是\emph{唯一可读的}。
一种被广泛使用的约定称为\emph{缩写项}，其规则如下：
% Part: lambda-calculus
% Chapter: syntax
% Section: abbreviated-syntax

\begin{enumerate}
\item 省略括号时，应用按从左到右进行。例如，若 $M$、$N$、$P$ 和~$Q$ 是项，则
  $MNPQ$ 是 $(((MN)P)Q)$ 的缩写。
\item 同样，省略括号时，λ 抽象被赋予尽可能宽的辖域。例如，$\lambd[x][MNP]$ 读作
  $(\lambd[x][MNP])$。
\item 单个 λ 可用于抽象多个变元。例如，$\lambd[xyz][M]$ 是
  $\lambd[x][\lambd[y][\lambd[z][M]]]$ 的缩写。
\end{enumerate}

例如，
\[
\lambd[xy][xxyx \lambd[z][xz]]
\]
是
\[
(\lambd[x][(\lambd[y][((((xx)y)x)(\lambd[z][(xz)]))])]).
\]
的缩写。

\begin{prob}
展开缩写项 $\lambd[g][(\lambd[x][g (x x)])
  \lambd[x][g (x x)]]$。
\end{prob}

\end{document}